\documentclass{article}

\usepackage{amsmath, amsthm, amssymb, amsfonts}
\usepackage{thmtools}
\usepackage{graphicx}
\usepackage{setspace}
\usepackage{geometry}
\usepackage{float}
\usepackage{hyperref}
\usepackage[utf8]{inputenc}
\usepackage[english]{babel}
\usepackage{framed}
\usepackage[dvipsnames]{xcolor}
\usepackage{tcolorbox}
\usepackage{physics}
\usepackage[shortlabels]{enumitem}

\DeclareMathOperator{\grd}{grad}
\DeclareMathOperator{\crl}{curl}
\DeclareMathOperator{\dvr}{div}
\DeclareMathOperator{\length}{length}
\DeclareMathOperator{\vspan}{span}
\DeclareMathOperator{\rng}{range}
\DeclareMathOperator{\nul}{null}
\colorlet{LightGray}{White!90!Periwinkle}
\colorlet{LightOrange}{Orange!15}
\colorlet{LightGreen}{Green!15}

\newcommand{\HRule}[1]{\rule{\linewidth}{#1}}

\declaretheoremstyle[name=Theorem,]{thmsty}
\declaretheorem[style=thmsty,numberwithin=section]{theorem}
\tcolorboxenvironment{theorem}{colback=LightGray}

\declaretheoremstyle[name=Proposition,]{prosty}
\declaretheorem[style=prosty,numberlike=theorem]{proposition}
\tcolorboxenvironment{proposition}{colback=LightOrange}

\declaretheoremstyle[name=Principle,]{prcpsty}
\declaretheorem[style=prcpsty,numberlike=theorem]{principle}
\tcolorboxenvironment{principle}{colback=LightGreen}

\setstretch{1.2}
\geometry{
    textheight=9in,
    textwidth=6in,
    top=0.5in,
    headheight=12pt,
    headsep=25pt,
    footskip=30pt
}

% ------------------------------------------------------------------------------

\begin{document}


% ------------------------'

\textbf{HW3}
All of the following exercises are from LADW.

\textbf{6.1.} True or false (explain your answers):

\begin{enumerate}
    \item[a)] Any system of linear equations has at least one solution;
    \item[b)] Any system of linear equations has at most one solution;
    \item[c)] Any homogeneous system of linear equations has at least one solution;
    \item[d)] Any system of $n$ linear equations in $n$ unknowns has at least one solution;
    \item[e)] Any system of $n$ linear equations in $n$ unknowns has at most one solution;
    \item[f)] If the homogeneous system corresponding to a given system of a linear equations has a solution, then the given system has a solution;
    \item[g)] If the coefficient matrix of a homogeneous system of $n$ linear equations in $n$ unknowns is invertible, then the system has no non-zero solution;
    \item[h)] The solution set of any system of $m$ equations in $n$ unknowns is a subspace in $\mathbb{R}^n$;
    \item[i)] The solution set of any homogeneous system of $m$ equations in $n$ unknowns is a subspace in $\mathbb{R}^n$.
\end{enumerate}

\textbf{7.1.} True or false (explain your answers):

\begin{enumerate}
    \item[a)] The rank of a matrix is equal to the number of its non-zero columns;
    \item[b)] The $m \times n$ zero matrix is the only $m \times n$ matrix having rank 0;
    \item[c)] Elementary row operations preserve rank;
    \item[d)] Elementary column operations do not necessarily preserve rank;
    \item[e)] The rank of a matrix is equal to the maximum number of linearly independent columns in the matrix;
    \item[f)] The rank of a matrix is equal to the maximum number of linearly independent rows in the matrix;
    \item[g)] The rank of an $n \times n$ matrix is at most $n$;
    \item[h)] An $n \times n$ matrix having rank $n$ is invertible.
\end{enumerate}


\end{document}